\section{الفصل~\ref{sec:glutcomputer}. ماذا تحسب عصبونات \nt{GLUT}}

القضايا المنطقية في الفصل مبرهنات عن دارات ذات أوزان غير سالبة، مستنتجة في الفصل نفسه؛ وتؤكد التجربة نموذج العصبون، وأن الدارات المبنية وفق هذه المبرهنات (كاشف التزامن، وخط التأخير، والمجموعة الذاتية الاستمرار، والسلسلة) موجودة فعلًا في الدماغ.

{\footnotesize
\setlength{\tabcolsep}{3pt}\renewcommand{\arraystretch}{1.15}
\begin{longtable}{>{\raggedright}p{0.5cm}>{\raggedright}p{4.0cm}>{\raggedright}p{5.6cm}>{\raggedright}p{2.3cm}>{\raggedright\arraybackslash}p{1.7cm}}
\toprule
م & القضية & التجربة وما قيس & المصدر & الحالة \\
\midrule\endfirsthead
\toprule
م & القضية & التجربة وما قيس & المصدر & الحالة \\
\midrule\endhead
1 & العصبون دارة RC بعتبة وتصفير~\eqref{eq:glutneuron} & حُقن في جسم العصبونات الهرمية في قشرة الجرذ (شرائح) تيار ضجيجي يشبه التدفق المشبكي في الدماغ الحي. علاقة التردد المتوسط بمتوسط التيار وتشتته يصفها العصبون المكامل التسريبي إذا أضيف إليه تكيّف بطيء في التردد. أما مجالات $\tau$ والتأخيرات فمذكورة في الفصل بلا إحالة إلى تجربة & \cite{Rauch2003} & مُقاس \\
2 & النموذج~\eqref{eq:glutneuron} تبسيط: فروع المداخل نفسها تطلق نبضات محلية & تسجيل مباشر من تغصنات العصبونات الهرمية في القشرة البصرية للفأر في الكائن الحي: المنبه البصري يستدعي نبضات تغصنية محلية تعتمد على مستقبلات \foreignlanguage{english}{NMDA}؛ وكبتها يوسّع ضبط العصبون على الاتجاه & \cite{Smith2013} & مُقاس \\
3 & نافذة التزامن $\Delta_{\max}=\tau\ln\frac{w}{\theta-w}$~\eqref{eq:window}؛ وعند $\theta=1.5w$ تساوي $0.7\tau$ & نتيجة للنموذج~\eqref{eq:glutneuron}. يوجد قياس مباشر لنافذة جمع ضيقة في عصبونات ذات تثبيط سريع (الفصل~\ref{sec:gaba}، السطر~3 من جدوله) & استنتاج في الفصل & نظرية \\
4 & شبكة من عصبونات \nt{GLUT} وحدها لا تحسب إلا دوال رتيبة في المداخل (القضية~\ref{prop:monotone}) & البرهان في الفصل: تكرار من الصمت بطرف أيمن غير متناقص. هذه قضية عن النموذج؛ ولا تجربة تختبرها على شبكة حية بلا تثبيط & استنتاج في الفصل & نظرية \\
5 & أعضاء الحس تعطي زوجًا من الأسلاك: بعض الخلايا يستجيب لاشتداد المنبه، وبعضها لضعفه & الشبكية: منذ الخلايا ثنائية القطب تتفرع إشارة المخاريط إلى قناتي تشغيل وإطفاء. الحجب الدوائي لخلايا التشغيل ثنائية القطب لدى القرد: القناتان تسيران منفصلتين حتى القشرة؛ وبعد الحجب يسوء اكتشاف اشتداد الضوء، لا خفوته & \cite{Schiller1992} & مُقاس \\
6 & بالشيفرة ثنائية السلك تحسب شبكة من عصبونات \nt{GLUT} أي دالة منطقية لا تزامنيًا، بلا إيقاع مشترك، بثمن مضاعفة عدد العصبونات (القضية~\ref{prop:dualrail}، \eqref{eq:dualrail}، الشكل~\ref{fig:dualxor}) & الشيفرة ثنائية السلك وتحديد الجاهزية من الإشارة نفسها تقنية معيارية في الدارات اللاتزامنية غير الحساسة للتأخير. ودارة \foreignlanguage{english}{XOR} من ستة عصبونات مبنية في الفصل. ولا تجربة تُظهر أن الدماغ يحسب الدوال المنطقية بالشيفرة ثنائية السلك تحديدًا & \cite{Sparso2001} & نظرية \\
7 & أكبر فرق في زمن وصول الصوت إلى الأذنين عند الإنسان $\approx0.6\ms$، والدماغ يميّز نحو عشرة أجزاء من المليون من الثانية & مستمعون في تجربة بخيارين للإجابة: عتبة تمييز فرق الزمن (75\,\% إجابات صحيحة) هي $9\,\mu\text{s}$ لضجيج في النطاق $150\text{--}1700\,\text{Hz}$، و$11\,\mu\text{s}$ لنغمة $1\,\text{kHz}$، و$28\,\mu\text{s}$ لنقرة. والحد $0.6\ms$ حساب الفصل، $0.22\,\text{m}/343\,\text{m/s}$ & \cite{KlumppEady1956} & مُقاس \\
8 & لدى الطيور يتحول فرق الأزمنة إلى موضع عبر خطوط التأخير وكواشف التزامن~\eqref{eq:itd} (الشكل~\ref{fig:itd}) & بومة الحظائر: تدخل المحاور الآتية من الأذنين نواة كواشف التزامن من جهتين متقابلتين؛ ويتغير زمن وصول النبضات على طول المحور بانتظام مع العمق، ومدى التأخيرات عبر النواة ($160\,\mu\text{s}$ على $700\um$) يطابق مدى التأخيرات بين الأذنين & \cite{Carr1990} & مُقاس \\
9 & لم يُعثر لدى الثدييات على صف من التأخيرات؛ وتحل المهمة نفسها كواشف تزامن بتثبيط محسوب التوقيت من عصبونات \nt{GLYC} & تسجيل في الكائن الحي في الزيتونة العلوية الإنسية لدى الجربوع: الاستجابات لا تشكل خريطة اتجاهات؛ وإيصال الغليسين وحاصره عبر ماصة دقيقة يُظهر أن التثبيط \emph{الغليسيني} المحسوب التوقيت بدقة يحدد مجال التأخيرات العامل. التثبيط هنا من \nt{GLYC} لا من \nt{GABA} & \cite{Brand2002,Grothe2010} & مُقاس \\
10 & المجموعة ذات التغذية الراجعة القوية قلّاب؛ والفعالية الذاتية الاستمرار تدوم ثوانيَ بعد المنبه (الشكل~\ref{fig:bistable}) & القشرة الجبهية الأمامية للقرد في مهمة حركة عين مؤجلة نحو نقطة محفوظة (تأخير $1\text{--}6\,\text{s}$): $87$ من $170$ عصبونًا مرتبطًا بالمهمة تغيّر ترددها أثناء التأخير، و$79\,\%$ منها يتوقف ذلك فيها على اتجاه النقطة المحفوظة. أما أن هذه الفعالية تحفظها تغذية راجعة بحالتين مستقرتين فنظرية & \cite{Funahashi1989,Wang2001} & مُقاس \\
11 & يُكتب البت إذا دام الوارد أكثر من $\approx0.7\tau$ (الشكل~\ref{fig:recurrentgroup}ب) & حل المعادلة $\tau\,\dd r/\dd t=-r+f(wr+I)$ بطريقة أويلر عند $w=1$ و$I=0.6$؛ ولا برنامج في \texttt{models/}. لا تجربة & حساب في الفصل & نموذج \\
12 & يمكن حفظ الذاكرة في تيسير المشابك، بلا نبضات تقريبًا & النظرية: الكالسيوم المتبقي في النهايات يحفظ المكتوب نحو ثانية. السند: في القشرة الجبهية الأمامية الإنسية لابن مقرض (تسجيل من عدة عصبونات معًا) توجد شبكة فرعية من العصبونات الهرمية متصلة بمشابك تيسيرية ذات أثر لاحق طويل. ومراجعة لأرصاد الفترات ”الصامتة“ أثناء الاحتفاظ بالذاكرة. أي الطريقتين تستعمل القشرة ومتى: مسألة مفتوحة & \cite{Mongillo2008,WangMarkram2006,Stokes2015} & فرضية \\
13 & تُمحى الذاكرة بالتعب: ينضب مخزون الناقل العصبي & تسجيلات مزدوجة لعصبونات هرمية في القشرة: تهبط استجابة المشبك مع النبضات المتواترة، وسرعة الهبوط يحددها احتمال القذف. أما أن هذا بالذات ما يطفئ المجموعة الذاتية الاستمرار فلم يُبيَّن مباشرة & \cite{Tsodyks1997} & مُقاس \\
14 & سلسلة المجموعات مؤقت يطلقه حدث؛ لدى الطيور المغردة تنطلق العصبونات بالتتابع الصارم & تسجيل عصبونات معرّفة في المركز الصوتي الأعلى لعصفور الزيبرا في النوم وأثناء الغناء: كل عصبون يرسل مخرجه إلى النواة الحركية يطلق رشقة قصيرة واحدة في لحظة واحدة دقيقة من الأغنية، وتنطلق العصبونات متتابعة & \cite{Hahnloser2002} & مُقاس \\
\bottomrule
\end{longtable}}
